\chapter{Version initiale de la page Signaux}

\section{Rôle de l'interface dans la chaîne méthodologique}
La page \texttt{/signals} constitue l'interface d'analyse détaillée de la couche
signal. Dans la version étudiée ici, elle prend en charge le passage entre le
résultat algorithmique produit par le pipeline A-G et l'interprétation humaine
nécessaire à la décision. Elle n'est donc ni une simple vue technique de
contrôle, ni un tableau de synthèse minimaliste : elle joue le rôle d'interface
explicative.

Sur le plan applicatif, cette page est structurée autour de
\texttt{frontend/app/signals/page.tsx},
\texttt{frontend/components/strategy/legacy-signals-view.tsx} et
\texttt{frontend/components/strategy/legacy-technical-analysis-panel.tsx}. Le
code conserve le terme \texttt{legacy} pour des raisons d'implémentation, mais la
présente étude parle de \emph{version initiale} ou de \emph{première version}
de la page Signaux.

\section{Organisation générale de l'écran}
L'écran se compose de deux ensembles. Le premier regroupe les contrôles de
contexte : sélection du symbole, choix de l'horizon et réglage éventuel du
\emph{cooldown}. Le second rassemble le contenu d'analyse. Cette séparation
permet de distinguer le paramétrage de la consultation du résultat.

Le choix de l'horizon occupe une place centrale dans cette interface. Un même
symbole peut en effet faire apparaître des lectures différentes selon que l'on
observe un régime de court terme, de moyen terme ou de long terme. La page ne se
contente donc pas d'afficher un signal ; elle rappelle implicitement que tout
résultat doit être rapporté à son horizon d'évaluation.

\section{Navigation par niveaux de lecture}
La première version de la page Signaux repose sur une navigation progressive en
plusieurs niveaux. Cette structuration n'est pas seulement ergonomique ; elle est
la traduction visuelle du pipeline méthodologique.

\subsection{Niveau 0 : vue d'ensemble agrégée}
Au premier niveau, l'utilisateur accède à un consensus général construit à partir
des familles effectivement disponibles. Cette vue répond à une question simple :
\emph{quelle est l'orientation globale du titre à l'horizon choisi ?}

Le score affiché à ce niveau est une synthèse, non un verdict opaque. Il sert de
porte d'entrée analytique. En pratique, il est calculé à partir des scores de
familles chargés pour le titre courant. Sa représentation sous forme de barre de
score permet une lecture immédiate de la polarité du consensus.

\subsection{Niveau 1 : lecture par catégories et par familles}
Le deuxième niveau expose les familles sous une organisation thématique :
tendance, momentum, oscillation et volume. Pour chaque famille, la page présente
un score, une étiquette interprétative et un éventuel état provisoire.

Ce niveau est fondamental du point de vue du rapport. Il matérialise la
contrainte d'explicabilité posée plus haut : un signal global ne doit jamais être
accepté comme une boîte noire. En affichant séparément les familles, la page
rend visible le fait qu'un consensus peut résulter d'une convergence forte entre
plusieurs lectures, ou au contraire d'une compensation partielle entre signaux de
natures différentes.

\subsection{Niveau 2 : inspection détaillée d'une famille}
Le troisième niveau est consacré à une famille donnée. Il fait apparaître
l'entonnoir de sélection du pipeline, depuis les variantes testées jusqu'aux
représentants retenus, ainsi que le score de famille et la liste des
représentants avec leurs poids.

Cette vue est particulièrement importante car elle relie l'interface à la
méthodologie. Le passage
\[
\text{testées} \rightarrow \text{viables} \rightarrow \text{compétitives} \rightarrow \text{représentatives}
\]
reprend exactement l'ordre des couches C, D et E. L'utilisateur voit donc non
seulement le résultat final, mais également le chemin de filtrage qui y conduit.

\section{Correspondance entre l'écran et le pipeline A-G}
La valeur de la page Signaux tient en grande partie à sa capacité à mettre en
scène le pipeline sans l'appauvrir. Cette correspondance peut être résumée ainsi.
\begin{itemize}
\item Les couches A à E sont visibles à travers les compteurs de l'entonnoir.
\item La couche F apparaît à travers les signaux courants des représentants.
\item La couche G apparaît à travers le score de famille et l'étiquette finale.
\end{itemize}

Autrement dit, l'écran ne montre pas seulement un résultat consolidé ; il expose
la logique de construction de ce résultat. Cette propriété rend l'interface plus
adaptée à un environnement professionnel qu'une simple liste de signaux non
justifiés.

\section{Gestion des états incomplets et honnêteté de restitution}
Une caractéristique importante de la page est la manière dont elle gère les cas
incomplets. Lorsqu'une famille n'est pas disponible, lorsqu'un résultat est
provisoire ou lorsqu'il n'existe pas suffisamment de données, l'interface le
signale explicitement.

Ce choix est loin d'être purement visuel. Il exprime une règle méthodologique :
un système de signalisation sérieux doit assumer les limites de ses résultats.
Forcer une orientation sur des bases insuffisantes reviendrait à dégrader la
qualité scientifique du pipeline lors de la restitution.

\section{Services applicatifs mobilisés}
La page interroge les services applicatifs dédiés à la couche signal, notamment
par les mécanismes associés à \texttt{useFamilyEnsemble(...)}. Le chargement des
familles est effectué en parallèle, ce qui permet une consultation fluide tout
en conservant une séparation stricte entre calcul et affichage.

Ce point est important dans le mémoire : la page Signaux n'effectue pas elle-même
la recherche quantitative. Elle consomme des résultats déjà produits selon le
pipeline défini dans le coeur applicatif, puis les restitue selon une logique de
lecture progressive.

\section{Apport fonctionnel de la page Signaux}
En définitive, la première version de la page Signaux apporte trois bénéfices
majeurs.
\begin{enumerate}
\item Elle permet une lecture immédiate du consensus global.
\item Elle rend possible une décomposition par familles cohérente avec la
structure économique du moteur de signaux.
\item Elle autorise un audit analytique jusqu'au niveau des représentants,
condition indispensable pour une utilisation crédible dans un contexte
professionnel.
\end{enumerate}

Ainsi, la page Signaux constitue l'interface privilégiée de compréhension du
moteur de signaux : elle ne remplace pas la méthode, mais elle la rend lisible,
contrôlable et exploitable.
